% Evans 9.7: printed pp.538--541; English PDF pp.555--558.
% Chinese: printed p.380 / scan p.398.
\section{Problems}
\label{sec:9-7}
The following short problem descriptions preserve the source numbers 1--14.
All solutions are editorial additions. Domain, decay, and regularity conditions
needed by a computation are stated explicitly.

\subsection*{1. A flow proof of the finite-dimensional zero lemma}
\begin{proof}[Solution]
Let $V$ be smooth near $\overline B_r$ and solve $x'=-V(x)$.
The inequality $V(x)\cdot x\geq0$ on the sphere makes the closed ball
forward invariant. One justification uses a Lipschitz bound near the ball:
for $d(x)=\operatorname{dist}(x,\overline B_r)$, projection to the sphere
gives $(d(x)^2)'\leq C d(x)^2$ outside, so an orbit starting with $d=0$
cannot leave by Gronwall. Boundedness in the compact ball permits global
continuation. For every $h>0$ its continuous time-$h$ map has a Brouwer
fixed point $y_h$. Take $h_j\downarrow0$ and $y_{h_j}\to y$.
Uniformly for initial points in the ball,
$x(h,y_h)=y_h-hV(y_h)+O(h^2)$, since $V$ and $DV$ are bounded there.
The fixed-point identity gives $V(y_h)=O(h)$, hence $V(y)=0$.
A fixed point of just one time map could be a periodic orbit; the
small-time limit is necessary to complete this argument.
\end{proof}

\subsection*{2. Weakly continuous scalar nonlinearities are affine}
\begin{proof}[Solution]
Write $G$ for the continuous scalar nonlinearity. Fix $s,t\in\R$,
$0<\theta<1$. On each interval of length $1/k$ let $f_k$ equal $s$ on
the first fraction $\theta$ and $t$ on the remainder.
For continuous test functions, cellwise averaging gives
$f_k\rightharpoonup m=\theta s+(1-\theta)t$; boundedness and density
extend this to all $L^2$ tests. The same argument gives
$G(f_k)\rightharpoonup\theta G(s)+(1-\theta)G(t)$. The assumed weak
continuity forces
$G(\theta s+(1-\theta)t)=\theta G(s)+(1-\theta)G(t)$.
A continuous function satisfying this identity on every segment is
$G(s)=\alpha s+\beta$: set $\beta=G(0)$, compare first on $[0,1]$,
then use segments containing $0,1$ and any real $s$. Conversely affine
maps preserve weak convergence by linearity.
\end{proof}

\subsection*{3. Penalty approximation of the nonnegative obstacle}
\begin{proof}[Solution]
Let $\beta_\varepsilon(s)=\min(s,0)/\varepsilon$ and solve
$-\Delta u_\varepsilon+\beta_\varepsilon(u_\varepsilon)=f$.
Existence follows by minimizing the strictly convex functional
\[
 \tfrac12\|Dv\|_2^2+{1\over2\varepsilon}\|v_-\|_2^2-\int fv
 \quad\hbox{on }H^1_0(U).
\]
Coercivity and convex lower semicontinuity give a minimizer, and
its derivative gives the equation. Testing by $u_\varepsilon$ bounds
$\|Du_\varepsilon\|_2$ and $\varepsilon^{-1}\|(u_\varepsilon)_-\|_2^2$.
Testing by $-(u_\varepsilon)_-$ gives more precisely
\[
 \|D(u_\varepsilon)_-\|_2^2+
 \varepsilon^{-1}\|(u_\varepsilon)_-\|_2^2
 \leq\|f_-\|_2\|(u_\varepsilon)_-\|_2,
\]
so $\|(u_\varepsilon)_-\|_2\leq\varepsilon\|f_-\|_2$.
A weakly convergent subsequence has limit $u\geq0$ by compact $L^2$
convergence. For each $w\geq0$ in $H^1_0$,
$\beta_\varepsilon(u_\varepsilon)(w-u_\varepsilon)\leq0$, hence
\[
 \int Du_\varepsilon\cdot D(w-u_\varepsilon)
 \geq\int f(w-u_\varepsilon).
\]
Rearrange with $\|Du_\varepsilon\|_2^2$ on the left and use weak
lower semicontinuity to obtain the desired variational inequality for $u$.
Two such solutions, tested against one another, give equality by strict
convexity. Thus every subsequential limit is $u$, giving convergence
of the full family weakly in $H^1_0$. Taking $w=u$ in the last inequality
also gives a matching upper bound for the gradient norm, so convergence
is in fact strong in $H^1_0$.
\end{proof}

\subsection*{4. Periodic forcing selects one periodic initial datum}
\begin{proof}[Solution]
For $f(t+\tau)=f(t)$, $f\in L^\infty$, define the period map
$P(g)=u(\tau;g)$ on $L^2(U)$. The difference of two solutions has zero
forcing, and its energy estimate gives
$\|P(g)-P(h)\|_2\leq e^{-\lambda_1\tau}\|g-h\|_2$.
Banach provides a unique $g_*$. The functions $u(t+\tau;g_*)$ and
$u(t;g_*)$ solve the same problem with the same initial trace; linear
uniqueness proves periodicity. Every periodic solution supplies a fixed
point, proving uniqueness of the initial datum as well.
\end{proof}

\subsection*{5. A small gradient nonlinearity}
\begin{proof}[Solution]
Assume $f\in L^2$ and $b$ is globally Lipschitz with constant $L$.
Let $Av$ solve $-\Delta Av=f-b(Dv)$, zero trace. Testing the difference
of two outputs gives
\[
 \|D(Av-Aw)\|_2^2\leq LC_P\|D(v-w)\|_2\|D(Av-Aw)\|_2.
\]
Thus $LC_P<1$ makes $A$ a contraction on $H^1_0$. Its fixed point
is unique and belongs to $H^2$ by the linear $L^2$ regularity theorem.
The value $b(0)$ affects the forcing bound but not the contraction factor.
\end{proof}

\subsection*{6. A positive solution above the first eigenvalue}
\begin{proof}[Solution]
Assume $U$ connected, $f$ bounded and Lipschitz, $f(0)=0$, and
$f'(0)>\lambda_1$. Normalize $\|\phi_1\|_\infty=1$.
For small $\varepsilon>0$, differentiability at $0$ gives
$f(\varepsilon\phi_1)\geq\lambda_1\varepsilon\phi_1$,
so $\underline u=\varepsilon\phi_1$ is a subsolution.
Let $-\Delta\psi=1$, $\psi=0$ on the boundary, and take
$\overline u=M\psi$ with $M\geq\sup|f|$.
It is a supersolution and dominates $\underline u$ after increasing $M$:
indeed $-\Delta[(\lambda_1\|\phi_1\|_\infty)\psi-\phi_1]\geq0$,
so comparison bounds $\phi_1$ by a multiple of $\psi$.
The ordered theorem yields $u\geq\varepsilon\phi_1>0$ in $U$.
This pointwise statement uses the continuous representative, supplied by
elliptic regularity with bounded right side. On a disconnected domain,
$f'(0)>\lambda_1(U)$ alone need not give positivity on every component;
connectedness is the needed qualification here.
\end{proof}

\subsection*{7. Classical comparison for the ordered iterates}
\begin{proof}[Solution]
With $h(s)=f(s)+\lambda s$ nondecreasing, each difference satisfies
$(-\Delta+\lambda)(u_{k+1}-u_k)=h(u_k)-h(u_{k-1})\geq0$
once $u_k\geq u_{k-1}$. Its boundary data are zero, except at the
first comparison where they are nonnegative by the subsolution trace.
The maximum principle starts and propagates the induction.
Likewise $(-\Delta+\lambda)(\overline u-u_{k+1})\geq
h(\overline u)-h(u_k)\geq0$. Thus the entire chain lies between the
barriers. The weak proof in \S9.3 is its positive-part counterpart.
\end{proof}

\subsection*{8. Critical bubbles and lost compactness}
\begin{proof}[Solution]
For $n\geq3$, $U(r)=(1+r^2)^{-(n-2)/2}$ has
\[
 \begin{split}
 U'&=-(n-2)r(1+r^2)^{-n/2},\\
 -U''-{n-1\over r}U'&=n(n-2)(1+r^2)^{-(n+2)/2}.
 \end{split}
\]
Thus $c=n(n-2)$. Scaling gives
$U_\lambda(x)=\lambda^{-(n-2)/2}U(x/\lambda)$ and the same equation.
For $2^*=2n/(n-2)$, changing variables $x=\lambda y$ proves invariance
of both $\|U_\lambda\|_{2^*}$ and $\|DU_\lambda\|_2$.
These norms are finite: at infinity the respective integrands have radial
orders $r^{-n-1}$ and $r^{1-n}$, and they are regular at the origin.
As $\lambda\downarrow0$, $U_\lambda(x)\to0$ for $x\ne0$.
Any strongly convergent $L^{2^*}$ subsequence would have limit zero
(by a further a.e. convergent subsequence), contrary to its constant
nonzero norm. This proves noncompactness along a specified sequence;
constant norms alone would not prove it.
\end{proof}
\Needspace{42mm}
\PDEAnnotation{集中需要极限证据}{尺度不变范数只给有界性。还须取 $\lambda\downarrow0$，证明几乎处处极限为零，再用强收敛保持范数得到矛盾。}

\subsection*{9. Energy and a weighted parabolic Pohozaev identity}
\begin{proof}[Solution]
Assume sufficient rapid decay for every displayed integral and boundary
flux, and $F'=f$, $F(0)=0$. Set $E=\int(|Du|^2/2-F(u))$.
Differentiate, integrate by parts, and use $u_t=\Delta u+f(u)$:
$E'=-\int u_t^2$. For the second moment $E_2$ of the same density,
\[
 E_2'=-\int |x|^2u_t^2-2\int u_t(x\cdot Du).
\]
The two cross terms are
$\int\Delta u(x\cdot Du)=(n-2)\int|Du|^2/2$ and
$\int f(u)(x\cdot Du)=-n\int F(u)$. Therefore
\[
 E_2'=-\int\bigl(|x|^2u_t^2-2nF(u)+(n-2)|Du|^2\bigr).
\]
The signs differ because the two densities have different scaling degrees.
Weighted integrability is essential; finite unweighted energy alone does
not license differentiating $E_2$.
\end{proof}

\subsection*{10. The normal cone and metric projection}
\begin{proof}[Solution]
For a nonempty closed convex $K\subset\R^n$, its indicator energy has
\[
 \partial I(x)=N_K(x)=\{v:v\cdot(y-x)\leq0\ (y\in K)\}
 \quad(x\in K),
\]
and empty subdifferential outside $K$. Minimizing
$I(x)+|x-z|^2/(2\lambda)$ gives $J_\lambda z=P_Kz$, independent
of $\lambda$, and $A_\lambda z=(z-P_Kz)/\lambda$.
The projection condition $z-P_Kz\in N_K(P_Kz)$ follows by varying
along each segment to $y\in K$. At an interior point the normal cone
is $\{0\}$; on a smooth boundary it is the outward normal ray.
\end{proof}

\subsection*{11. An irreversible flow}
\begin{proof}[Solution]
Take $H=\R$, $I(s)=|s|$. Then $\partial I(s)=\{\operatorname{sgn}s\}$
for $s\ne0$ and $\partial I(0)=[-1,1]$. Directly,
$S(t)s=\operatorname{sgn}(s)(|s|-t)_+$ solves the inclusion a.e.
For example $S(2)1=S(2)2=0$ although $1\ne2$. This semigroup is
contractive and continuous but its time maps need not be injective;
a backward inverse flow cannot be defined on all states.
\end{proof}

\subsection*{12. A constrained sandpile as a normal-cone flow}
\begin{proof}[Solution]
Let $K=\{v\in L^2(\R^2):|Dv|\leq1\text{ a.e.}\}$, with the
inequality understood in the weak-gradient sense. It is convex and closed
in $L^2$: distributional limits preserve the gradient bound, by weak-*
compactness of bounded gradients on compact subsets. Its indicator $I$
is proper, convex and lower semicontinuous.
Assume the proposed fields are regular and integrable enough that
$f-u_t=-\operatorname{div}(aDu)\in L^2$ and integration by parts
against $v-u$ is justified, for instance by compactly supported flux
and approximation. For $v\in K$,
\[
 (f-u_t,v-u)=\int aDu\cdot(Dv-Du)\leq0.
\]
Indeed where $a>0$ one has $|Du|=1$, so
$Du\cdot Dv-|Du|^2\leq|Dv|-1\leq0$; elsewhere the integrand is zero.
This is precisely $f-u_t\in N_K(u)=\partial I(u)$.
The calculation is conditional on the asserted flux and regularity;
it does not construct $a$ from a solution of the abstract inclusion.
\end{proof}

\subsection*{13. Exponential decay of the speed}
\begin{proof}[Solution]
For a smooth solution with compatible boundary data, put $v=u_t$.
Differentiating $u_t-\operatorname{div}DL(Du)=0$ gives
$v_t-\operatorname{div}(D^2L(Du)Dv)=0$, with $v=0$ on the boundary.
Testing by $v$, uniform convexity and Poincar\'e give
\[
 \tfrac12(\|v\|_2^2)'\leq-\theta\|Dv\|_2^2
 \leq-\theta\lambda_1\|v\|_2^2.
\]
Hence $\|u_t(t)\|_2^2\leq\|u_t(0)\|_2^2e^{-2\theta\lambda_1t}$
when the initial speed is finite. Otherwise the same estimate starts
at any positive regular time, with a constant depending on that time.
The uniform coefficient $\theta$ is the source of the decay rate.
\end{proof}

\subsection*{14. Total variation under nonlinear diffusion}
\begin{proof}[Solution]
Write $v=\phi(u)$ with $\phi'>0$. Under the source's nonvanishing-gradient
assumption and sufficient integrability and vanishing flux at infinity,
$Du/|Du|=Dv/|Dv|$. Thus
\[
 {d\over dt}\int|Du|=\int{Dv\over|Dv|}\cdot D\Delta v.
\]
Differentiating $|Dv|$ twice gives the pointwise Bochner identity
\[
 \Delta|Dv|={Dv\over|Dv|}\cdot D\Delta v
 +{|D^2v|^2\over|Dv|}-{|D^2vDv|^2\over|Dv|^3}.
\]
Its Laplacian integrates to zero under the stated flux hypothesis, so
\[
 {d\over dt}\int|Du|
 =-\int\left({|D^2v|^2\over|Dv|}
       -{|D^2vDv|^2\over|Dv|^3}\right)\leq0.
\]
The inequality uses $|D^2vDv|\leq|D^2v||Dv|$.
This is the requested smooth calculation where denominators are defined.
At critical points one needs a regularization and a limiting argument;
no general BV contraction theorem is proved here, and nonvanishing and
whole-space decay should not be silently imposed on arbitrary data.
\end{proof}
